\begin{table}[htbp]\centering
\caption{Measurement-invariance screen: regression-based DIF, IT vs.\ IA/Cyber.}\label{tab:invariance}
%% ALT: Measurement-invariance screen using regression-based differential item functioning (DIF) between IT and IA/Cyber respondents: per-item beta, uniform-DIF p-value, and flag status.
\begin{tabular}{@{}lrrc@{}}\toprule
\textbf{Item} & \textbf{$\beta_{\text{IA}}$ (matched)} & \textbf{$p$} & \textbf{DIF flag} \\ \midrule
Q1.1 & -0.99 & $<.001$ & yes \\
Q2.5 & +0.23 & $<.001$ & yes \\
Q5.2 & +0.16 & .003 & no \\
Q2.4 & +0.16 & $<.001$ & no \\
Q2.6 & +0.16 & $<.001$ & no \\
Q2.1 & +0.15 & $<.001$ & no \\
Q2.2 & +0.15 & $<.001$ & no \\
Q2.3 & +0.14 & $<.001$ & no \\
\bottomrule
\end{tabular}

\smallskip
\parbox{4.8in}{\footnotesize \textit{Note.} Post-hoc analysis, not specified in the proposal. $\beta_{\text{IA}}$ is the
IA-vs-IT difference in item response (scale points) after matching on the item's composite score
(a linear implementation of Zumbo's DIF framework). 2 of 23 items met the flag criterion ($p<.001$ and
$|\beta_{\text{IA}}| \geq 0.20$). All flagged effects are small relative to the composite-level
field difference, which approaches two standard deviations.}
\end{table}